\section{실험 설정}\label{sec:setup}
\subsection{학습 자료와 평가 집합}
추가 보정기에는 합성 코너 감독을 사용한다. 등록된 합성 자료 60,000장은 학습 55,980장, 추론 온도 선택용 calibration 1,004장, selection 1,031장, held-out 1,985장으로 나뉜다. 뒤의 두 분할은 원래 등록된 역할을 보존하되, 본문의 실사 결과나 독립 평가의 근거로 사용하지 않는다. 이 두 분할의 현재 실험 내 추가 사용 이력은 확인이 필요하다. 초기 추정기의 일반 영상 사전학습과 추가 보정기 학습은 구분하며, 전체 시스템의 모든 사전학습이 합성 전용이었다고 가정하지 않는다.

주 실사 집합은 13개 촬영 세션의 직사각형 영상 319장으로, 플라스틱 194장과 목재 125장을 포함한다. 이 집합은 개발과 분석에 반복 사용되었으며 독립 미사용 시험집합이 아니다. 서로 다른 실제 팔레트 개체 수, 프레임 모집·제외의 전체 이력 및 카메라별 적용 범위는 아직 완전하게 확인되지 않았다. 따라서 등록된 영상 집합에서의 효과로 해석하며 실제 운용 빈도에 대한 대표성을 주장하지 않는다.

319장의 가능한 코너 자리는 $319\times8=2,552$개이고, 평가 참조가 있는 코너는 2,499개이다. 참조가 없는 53자리와, 참조는 있지만 YOLO의 조건부 평가에 포함되지 않은 54코너를 구분한다. YOLO의 객체 매칭 영상은 311장, 유효 평가 예측은 2,445코너이다. 참조 좌표에는 수동 표시와 기하 보완이 포함되며, 각 생성 경로의 정확한 비율은 현 집합에서 추가 확인이 필요하다.

정사각형 GREEN0918은 등록 치수 $[1.1,1.1,0.15]$m인 팔레트의 119장·한 촬영 세션이다. 동일한 고정 모델에 대해 선언된 수동 코너 602개와 화면 밖 좌표 2개를 제외한 600개 모드를 각각 보고한다. 두 모드는 동일 영상의 참조 포함 규칙만 다르다. 가능한 952자리 중 350자리는 전체 모드에도 참조가 없으며, 이를 예측 실패로 세거나 전부 가려진 코너로 간주하지 않는다. 과거 실사 지도자료와의 개체·촬영·학습 이력 전체 대응은 미확인이므로 미관측 객체 일반화로 부르지 않는다. 독립 표준 자세 참조가 없어 정사각형 위치·회전 정확도는 \notrun 으로 남긴다.

\begin{table}[t]
\centering\footnotesize
\setlength{\tabcolsep}{4pt}
\caption{학습·평가 자료의 역할과 확인 범위. 실사 집합들은 하나의 독립 시험집합으로 합산하지 않는다.}\label{tab:data}
\begin{tabularx}{\columnwidth}{p{.28\columnwidth}rY}
\toprule
자료 & 영상 수 & 역할 \\
\midrule
합성 학습 & 55,980 & 보정기 합성 감독 \\
합성 calibration & 1,004 & 추론 온도 선택 \\
합성 selection & 1,031 & 등록 분할; 추가 사용 이력 확인 필요 \\
합성 held-out & 1,985 & 등록 분할; 별도 평가 완료를 뜻하지 않음 \\
\midrule
직사각형 실사 & 319 & 13세션; 재사용 개발 평가 \\
플라스틱 / 목재 & 194 / 125 & 319장의 하위집단 \\
정사각형 실사 & 119 & 1세션; 별도 2D 평가 \\
업데이트 대안 & 128 & 공통 비노출 개발 부분집합 \\
\bottomrule
\end{tabularx}
\tabnote{합성 네 분할 합계는 60,000장이다. 초기 추정기의 사전학습과 보정기의 추가 학습을 구분한다. 두 128장 집합은 frame ID가 전수 일치하며, 개체 수와 모집·제외 이력은 추가 확인 대상이다.}
\end{table}


\subsection{사람 입력 등급과 코너 가시성}
표~\ref{tab:composition}은 원본 RGB 해시와 frame ID에 연결한 최신 사람 입력 등급의 구성이다. 직사각형 319장은 없음 153장, 중간 92장, 어려움 74장이며 플라스틱은 72/57/65장, 목재는 81/35/9장이다. 정사각형119도 없음 3장·중간 85장·어려움 31장으로 입력되었다. 이는 등급 입력 완료를 뜻하며 판정 기준의 사람 확인과는 구분된다. 현재 판정 기준은 미확인 상태이므로 외부 가림 검수 완료로 표현하지 않고 보조 사후 분석으로 제시한다.

외부 물체 가림·팔레트 자체 가림·화면 밖 잘림은 서로 다른 속성이다. 최신 사람 입력의 중간·어려움 경계와 이 속성의 포함 여부를 CLI가 추정해 확정하지 않았다. 과거 외부 가림 128장 패널은 기존29/20/79와 더 오래된29/21/78 버전 차이를 포함한 원래 출처와 숫자로 별도 보존하며 최신319장 등급과 같은 계약으로 합치지 않는다. 임의로 개수를 맞추거나 오차에 따라 집단을 재선택하지 않았다.

영상 단위 등급과 코너 단위 가시성도 구분한다. 기존 잠금71점과 새 입력3030점을 중복 없이 연결해 직사각형2499점·정사각형602점의 총3101개 참조 코너 상태를 보존하였다. 오차 계산에는 원래 평가 참조 좌표와 같은 전체 객체 대응을 사용하며, 새로운 재클릭이나 PnP 생성점으로 정답 좌표를 교체하지 않았다. 가시성 입력과 좌표의 생성 경로·독립 블라인드 검수 여부는 서로 다른 정보이다. 상세 결과와 한계는 보충자료에 둔다.
\begin{table}[t]
\centering\footnotesize
\setlength{\tabcolsep}{3pt}
\caption{최신 사람 입력 등급(기준 미확인)의 구성.}\label{tab:composition}
\begin{tabularx}{\columnwidth}{Yrrrrr}
\toprule
집단 & 없음 & 중간 & 어려움 & 미입력 & 전체 \\
\midrule
직사각형 / 플라스틱 & 72 & 57 & 65 & 0 & 194 \\
직사각형 / 목재 & 81 & 35 & 9 & 0 & 125 \\
직사각형 합계 & 153 & 92 & 74 & 0 & 319 \\
정사각형119 & 3 & 85 & 31 & 0 & 119 \\
\bottomrule
\end{tabularx}
\tabnote{숫자는 실제 입력된 등급의 개수이며 판정 기준 승인과 구분한다. 과거 외부 가림 128장 레이블과 혼합하지 않는다. 과거 정사각형 150장은 별도 7세션 자료로 현 119장 성능 분모에 넣지 않았다.}
\end{table}


\subsection{비교 방법과 추가 학습}
Base는 같은 기반 추정기의 \textbf{보정 없는 출력}이다. 세 기반은 YOLO, DOPE(Deep Object Pose Estimation)\cite{dope2018}, ResNet(Residual Network)-18 기반 팔레트 9점 모델\cite{simple2018}이다. 같은 Base 출력에 N3를 적용한 전후를 비교하며, 특징 연결층과 보정기는 기반별로 학습한다. 이 실험의 목적은 기반끼리의 절대 순위가 아니라 동일 보정 절차의 적용성이다.

구성 요소 비교는 YOLO에서 영상 기반 국소 보정 N0, 회전 대응 감독만 추가한 N1, 치수만 추가한 N2, 둘을 결합한 N3를 사용한다(표~\ref{tab:ablation_contract}). 초기화·표본 순서·추가 학습량을 맞춘 비교이지만 치수 분기의 정보 내용과 추가 파라미터 효과를 완전히 분리하지는 못한다. 기존 독립 실행의 영상 보정 P는 N0와 다른 가중치이며, 전체 방법 비교와 보충자료에서 별도로 구분한다.
\begin{table}[t]
\centering\footnotesize
\setlength{\tabcolsep}{4pt}
\caption{구성 요소 대조. 모든 경로의 PnP에는 같은 치수를 사용한다.}\label{tab:ablation_contract}
\begin{tabularx}{\columnwidth}{lYcc}
\toprule
구성 & 보정 & \shortstack{보정기의\\치수 입력} & \shortstack{학습 시\\회전 대응} \\
\midrule
Base & 없음 & -- & -- \\
N0 & 영상 기반 국소 보정 & 없음 & 고정 번호 \\
N1 & 영상 기반 국소 보정 & 없음 & 동등성 고려 \\
N2 & 영상 기반 국소 보정 & 사용 & 고정 번호 \\
N3 & 영상 기반 국소 보정 & 사용 & 동등성 고려 \\
\bottomrule
\end{tabularx}
\tabnote{N0--N3가 구성 요소 비교이다. 별도 실행의 영상 보정 P를 N0로 재명명하지 않는다.}
\end{table}


N3는 각 기반에서 3개 seed로 학습한다. seed는 초기화와 표본 순서에 사용하는 난수 설정이다. 각 학습은 batch 16·6,000회 업데이트·96,000회 표본 노출이며, AdamW의 기준 학습률 $10^{-3}$, 가중치 감쇠 $10^{-4}$, 100회 준비 구간과 코사인 감소, 최종 학습률 0.1배, 기울기 노름 상한 5를 사용한다. 6,000회째 가중치를 평가한다. 유효 고유 학습 행은 DOPE 44,063개·ResNet 55,806개로, 같은 노출량이 같은 고유 자료 수나 기본 모델까지의 동일 총학습량을 뜻하지 않는다.

학습 손실의 온도는 1이다. 추론 온도는 합성 calibration에서 $\{0.5,1,2,4\}$ 중 고정 번호 목표의 교차엔트로피로 선택하며 실사 성능으로 재선택하지 않는다. DOPE와 ResNet의 여섯 실행은 모두 온도 1을 선택하였다. 후보 구성과 원본 이동 상한은 평가 전에 고정한다.

전체 방법 비교에는 직접 회귀 D, 선 구조 보정 L, PoseFix 기반 팔레트 보정과 N3를 포함한다. PoseFix 비교 구현은 공식 네트워크를 기준으로 연산을 대조한 PyTorch 팔레트 9점 변형이며 입력·목표 및 일부 수치 설정을 변경하였다. 3개 seed에서 각각 6,000회 업데이트·96,000회 표본 노출의 최종 가중치를 사용한다. 이는 원 논문의 전체 수렴 학습을 재현했다는 뜻은 아니다. 입력·출력·손실·감독 조건의 차이를 가진 전체 방법 비교이며, D/P의 차이도 분포 출력 하나의 효과로 해석하지 않는다. 비교군의 세부 사전학습·설정 선택과 비용의 완전한 동등성은 미확인 부분으로 남긴다.

추정기 업데이트의 대안은 별도 플라스틱 128장에 대해 비교한다. 원래 학생·보정 의사 레이블 학생·합성 추가 학습·Base·Base+N3를 모두 같은 영상에서 평가한다. 이 집합은 공통 비노출 조건이 확인된 재사용 개발 부분집합이며 독립 시험은 아니다. frame ID 전수 대조에서 기존 가림 128장과 교집합 128장, 양쪽 차집합 0장으로 확인되었다.

\subsection{코너 오차, 자세 오차와 실패 분모}
코너 정확도는 원본 영상에서 예측 코너와 참조 코너 사이의 유클리드 거리로 평가한다. 회전 동등성을 허용할 때는 하나의 전체 객체 순열을 적용하고 점별 자유 최근접 매칭은 하지 않는다. \emph{코너 오차 중앙값·P90(px)}는 객체 매칭을 통과한 유효 예측 코너에 대한 조건부 통계이다. P90은 90번째 백분위수이며 상위 10\%의 평균이 아니다. 코너 거리 기반 평가는 팔레트 키포인트의 직접적인 정확도를 다루는 선행\cite{mueller2024}과 같은 평가 목적이지만, 평균·중앙값·축척이 다르면 보고 숫자를 직접 비교하지 않는다.

Percentage of Correct Keypoints(PCK)는 임계 거리 이내의 코너 비율이다. \pckten 은 10픽셀 이내인 예측 수를 \textbf{전체 주석 코너 수}로 나눈 백분율이며 결측·매칭 실패는 정답으로 인정하지 않는다. 주석이 없는 물리적 코너까지 평가했다는 의미는 아니다. 전체 실패 벌점 P90은 결측·매칭 실패에 영상 대각선 길이를 부여한 별도 통계이다. 조건부 오차와 이 전체 통계를 함께 해석한다.

위치 오차 $e_t$는 미터 단위 위치 벡터의 거리, 회전 오차 $e_R$는 객체가 허용하는 회전 집합 $\mathcal G$ 안의 최소 회전각 차이이다. 참조는 $\R^*,\tvec^*$, 예측은 $\hat\R,\hat\tvec$로 표시한다.
\begin{align}
 e_t&=100\|\hat\tvec-\tvec^*\|_2\quad[\mathrm{cm}],\\
 e_R&=\frac{180}{\pi}\min_{Q\in\mathcal G}\arccos\!\left[\operatorname{clip}\!\left(\frac{\operatorname{tr}((\R^*Q)^\top\hat\R)-1}{2},-1,1\right)\right].
\end{align}
객체 대칭을 명시적으로 정의한다는 평가 원칙은 BOP\cite{bop2018}와 관련되지만, 각 프로젝트의 점·참조 정의가 같은 것은 아니다. 본 위치·회전 참조는 영상 주석과 물체 기하로 재구성한 값이며 독립 물리 계측의 정답이 아니다. 자세가 산출된 영상의 중앙값·P90과 전체 영상에 대한 자세 산출률을 함께 보고한다. 산출률은 정답률이 아니며, 정답 박스 매칭을 자세 추론의 입력 또는 성공 게이트로 사용하지 않는다. 따라서 코너 매칭 영상과 자세 산출 영상의 집합은 같지 않을 수 있다.

방향각 오차, 방향을 고려한 3차원 Intersection over Union(IoU), 8코너에 전체 회전 동등성을 적용한 Average Distance of Model Points(ADD)의 변형 $\mathrm{ADD}_{\rm sym}$은 보충자료에 제시한다. 정확도 곡선 아래 면적 AUC(Area Under the Curve)는 객체 대각선으로 정규화한 오차의 0--0.1 범위에서 적분·정규화하며 자세 실패를 전체 분모에 남긴다. 이는 자유로운 최근접 표면 대응 ADD-S와 다르다. 추가 영상 지표의 정의와 원수치도 보충자료에 보존한다.

\subsection{집계, 불확실성과 큰 오류 분석}
학습 방법의 값은 각 seed에서 계산한 통계량의 산술평균이다. 예측 앙상블이나 모든 seed의 점을 한꺼번에 모은 중앙값이 아니다. 13개 세션을 단위로 10,000회 짝지어 재표집하고, 각 재표집 안에서 세 seed의 통계 차이를 평균해 2.5·97.5백분위 구간을 계산한다. 초기 난수는 20260917이며 같은 재표집을 사용한다. 주 변화량은 $\operatorname{median}(e_{\rm after})-\operatorname{median}(e_{\rm before})$로, 프레임별 차이의 중앙값과 구분한다. 이 구간은 반복 개발자료의 사후 기술 통계이며 다중 비교·개발 중 선택 편향을 보정한 독립 확증이 아니다. 정사각형 한 세션의 세션 간 일반화 구간은 산출하지 않는다.

큰 오류 분석은 초기 예측에서 정한 하나의 참조 대응을 보정 후에도 고정한다. 주 코너 지표가 각 예측의 허용 전체 대응을 선택하는 것과 구분하여 이동 상한의 효과를 진단한다. 초기 오차가 한계 안인지 밖인지, 원래 좋은 코너의 손상, 큰 오류 복구, 2D와 위치·회전의 공동 변화를 기록하며 표본 제외나 파라미터 선택에는 사용하지 않는다. 미측정·미확인 수치는 \notrun, 정의되지 않는 통계는 NA로 표기한다.
